\chapter{Experimental results}
\label{ch:results}%
% Aim: 20-25 pages. Present observations first; interpretation comes at the end of the chapter and in Chapter 5.

\section{Experimental programme}
\label{sec:programme}
% Table of all tests: ID, ice content, temperature, target pressure, corrected stress,
% initial height and density, duration, remarks. State which tests are excluded and why.

\section{Pure ice}
\label{sec:results_ice}
% Validation of procedure and data reduction against the published pure-ice compaction data from the same vessel.
\subsection{Compaction during loading}
\subsection{Compaction under constant stress}
% Height versus log-time; densities at fixed times; strain rate versus density.

\section{Ice--regolith mixtures}
\label{sec:results_mixtures}
\subsection{Effect of temperature}
\subsection{Effect of ice content}
\subsection{Effect of axial stress}

\section{Dry regolith}
\label{sec:results_dry}
% Reference test for the skeleton-only end-member; comparison with literature oedometric data.

\section{Unloading and reloading}
\label{sec:results_unloading}
% Dead band due to seal friction; unloading modulus from the slope where the piston moves.

\section{Constant-rate compression and stress relaxation}
\label{sec:results_relaxation}
% Pressure versus piston travel at constant rate; pressure decay in log-time at locked oil volume.
% State clearly what can and cannot be concluded from these tests.

\section{Post-test observations}
\label{sec:results_posttest}
% Final height and mass, mass gained by the frits, visual appearance, microstructure if available.

\section{Summary of experimental findings}
\label{sec:results_summary}
